\documentclass[conference]{IEEEtran}

% --- Language and Encoding ---
\usepackage[T1]{fontenc}        % Suporte para acentuação e fontes de 8 bits
\usepackage[utf8]{inputenc}     % Codificação de caracteres moderna
\usepackage[brazil]{babel}      % Tradução de termos e hifenização em português

% --- TikZ Setup (For Diagrams) ---
\usepackage{tikz}
\usetikzlibrary{shapes.geometric, arrows.meta}

% Configurando estilos do TikZ para a árvore de comportamento
\tikzset{
    control/.style={rectangle, draw, fill=gray!15, minimum width=2.8cm, minimum height=0.8cm, font=\normalsize\bfseries},
    action/.style={rectangle, draw, fill=blue!10, rounded corners=3pt, minimum width=2.8cm, minimum height=0.8cm, font=\scriptsize\ttfamily},
    condition/.style={rectangle, draw, fill=green!10, rounded corners=3pt, minimum width=2.8cm, minimum height=0.8cm, font=\scriptsize\ttfamily}
}

% --- Algorithm Setup (For Pseudocode) ---
\usepackage{amsmath}
\usepackage[linesnumbered,lined,titlenumbered,ruled,noend]{algorithm2e}

% Redefine comment delimiters to match the reference: /* ... */
\SetCommentSty{texttt}
\SetKwComment{Comment}{/* }{*/}

\begin{document}

% --- Title and Authors ---
\title{Título do Artigo Científico para a IEEE}

\author{
    \IEEEauthorblockN{Nome do Autor 1}
    \IEEEauthorblockA{Departamento de Computação\\
    Universidade Exemplo (UE)\\
    Cidade, Estado\\
    Email: autor1@exemplo.com}
    \and
    \IEEEauthorblockN{Nome do Autor 2}
    \IEEEauthorblockA{Departamento de Elétrica\\
    Universidade Exemplo (UE)\\
    Cidade, Estado\\
    Email: autor2@exemplo.com}
}

\maketitle

% --- Abstract ---
\begin{abstract}
Este resumo deve descrever brevemente o objetivo do seu artigo da IEEE, a metodologia utilizada, os resultados obtidos e as conclusões principais. O documento foi estruturado especificamente para aceitar a escrita nativa em língua portuguesa e renderizar gráficos utilizando a biblioteca TikZ.
\end{abstract}

\begin{IEEEkeywords}
Computação, LaTeX, Diagramas TikZ, IEEEtran.
\end{IEEEkeywords}

% --- Sections ---
\section{Introdução}
Esta é a seção de introdução. Note que a acentuação funciona nativamente (por exemplo: acentuação, coração, computação) graças aos pacotes configurados. 

\section{Metodologia}
Abaixo, apresentamos o diagrama da árvore de comportamento do sistema, modelada no Groot 2 e representada em vetor pelo TikZ.

\begin{figure}[!ht]
    \centering
    % Início do ambiente da árvore TikZ
    \begin{tikzpicture}[
        sibling distance=3cm,
        level distance=1.5cm,
        every node/.style={align=center},
        edge from parent/.style={draw, -Latex, thick}
    ]
        \node [control] {?}
            child {node [control] {$\rightarrow$}
                child {node [condition] {New position \\ requested}}
                child {node [control] {?}
                    child {node [control] {$\rightarrow$}
                        child {node [action] {GNSS \\ fix}}
                        child {node [action] {Send \\ position}}
                    }
                    child {node [action] {No position \\ alert}}
                }
            }
            child {node [action] {Sleep}};
    \end{tikzpicture}
    \caption{Árvore de comportamento equivalente obtida do arquivo XML.}
    \label{fig:diagrama_exemplo}
\end{figure}

Como observado na Fig. \ref{fig:diagrama_exemplo}, a árvore de comportamento é compilada sem distorções no PDF final.

\subsection{Behavior Tree Tick Algorithm}
The tick process traverses the tree iteratively using a DFS approach over the LCRS (Left-Child Right-Sibling) representation, as shown in Algorithm~\ref{alg:btf_tick_tree}.

\begin{algorithm}[!ht]
\caption{Behavior Tree Tick}
\label{alg:btf_tick_tree}
\SetAlgoLined
\SetKwInOut{Input}{Input}
\SetKwInOut{Output}{Output}
\SetKwFunction{FRoot}{root}
\SetKwFunction{FAction}{action}
\SetKwFunction{FFirstChild}{first\_child}
\SetKwFunction{FControl}{control}
\SetKwFunction{FNextSib}{next\_sibling}
\SetKw{KwAnd}{and}
\SetKw{KwNot}{not}
\Input{Behavior tree with LCRS node representation}
\Output{Final status of the root node}
$node \leftarrow \FRoot{tree}$\;
$status \leftarrow \text{UNDEFINED}$\;
\While{\KwNot ($status$ is defined \KwAnd $node$ is root)}{
    \eIf{$status = \text{UNDEFINED}$}{
        \eIf{$node$ is a leaf}{
            $status \leftarrow node.\FAction{tree}$\;
            \If{$status = \text{RUNNING}$}{
                store $node$ as $running\_node$\;
            }
        }{
            $node \leftarrow \FFirstChild{node}$\;
        }
    }{
        $parent \leftarrow node.parent$\;
        $result \leftarrow parent.\FControl{tree, node, status}$\;
        \eIf{$result = \text{CONTINUE}$}{
            \eIf{$node$ has sibling}{
                $node \leftarrow \FNextSib{node}$\;
                $status \leftarrow \text{UNDEFINED}$\;
            }{
                $node.parent.status \leftarrow status$\;
                $node \leftarrow node.parent$\;
            }
        }{
            $node.parent.status \leftarrow status$\;
            $node \leftarrow node.parent$\;
        }
    }
}
\Return{$status$}\;
\end{algorithm}

\section{Conclus\~ao}
Seu ambiente está pronto para produzir documentos acadêmicos robustos e compatíveis com as normas da IEEE.

\end{document}